% -----------------------------------------------------------
\section{Word}
% -----------------------------------------------------------
	\label{WORD}
	
% % ----------------------
% \subsection{See also}
% % ----------------------

% % ----------------------
% \subsection{1828 American Dictionary of the English Language} % *1828
% % ----------------------
	% \label{WORD-1828}

	% \begin{compactitem}
		% \item
	% \end{compactitem}

% % ----------------------
% \subsection{Oxford English Dictionary} % *OED
% % ----------------------
	% \label{WORD-OED}

	% \begin{compactitem}
		% \item[]
	% \end{compactitem}

% ----------------------
\subsection{Scriptural Usage} % *SCRIPTURES
% ----------------------
	\label{WORD-SCRIPTURES}

	% % -----
	% \subsubsection{Old Testament} % *OT
	% % -----
		% \label{WORD-OT}
	
		% % --
		% \paragraph{KJV}
		% % --
			% \label{WORD-OT-KJV}
	
			% \begin{tabular}{ll}
				% Total occurrences in the OT & 
			% \end{tabular}
			
			% \begin{compactdesc}
				% \item[]
			% \end{compactdesc}
			
		% % --
		% \paragraph{Hebrew Bible}
		% % --
			% \label{WORD-HEBREW}
		
			% %
			% \subparagraph{\texorpdfstring{\he{sample word}}{}} % paste actual hebrew text here
			% %
				% \label{PAGA} % whatever is in step bible without periods
			
				% \begin{compactdesc}
					% \item[HALOT , PDF ] \textbf{} % Use the 1994 PDF
						% \begin{compactitem}
							% \item[1]
						% \end{compactitem}
					% \item[BDB , PDF ] \textbf{}
						% \begin{compactitem}
							% \item[1]
						% \end{compactitem}
					% \item[DCH PDF ] \textbf{} % don't include the actual page number, they repeat from volume to volume
						% \begin{compactitem}
							% \item[1]
						% \end{compactitem}
				% \end{compactdesc}
				
				% \begin{compactdesc}
					% \item[Some OT uses] \textbf{}
						% \begin{compactdesc}
							% \item[]
						% \end{compactdesc}
				% \end{compactdesc}
			
		% % --
		% \paragraph{Septuagint}
		% % --
			% \label{WORD-LXX}
		
			% %
			% \subparagraph{\texorpdfstring{\gr{sample word}}{}}
			% %
		
	% -----
	\subsubsection{New Testament} % *NT
	% -----
		\label{WORD-NT}
	
		% % --
		% \paragraph{KJV}
		% % --
			% \label{WORD-NT-KJV}		
	
			% \begin{tabular}{ll}
				% Total occurrences in the NT & 
			% \end{tabular}
			
			% \begin{compactdesc}
				% \item[]
			% \end{compactdesc}
			
		% --
		\paragraph{Greek New Testament} % *GREEK
		% --
			\label{WORD-GREEK}
		
			%
			\subparagraph{\texorpdfstring{\gr{λόγος}}{}}
			%
				\label{LOGOS}
			
				\begin{compactdesc}
					\item[BDAG 530b] \textbf{}
						\begin{compactitem}
							\item[1] \textbf{a communication whereby the mind finds expression, \textit{word}}
							\item[1.a] of utterance, chiefly oral
							\item[1.a.$\alpha$] as expression, \textit{word} (oratorical ability plus exceptional performance were distinguishing marks in Hellenic society, hence the frequent association of \gr{λ}. and \gr{ἔργον} `deed'...
							\item[1.a.$\beta$] The expression may take on a variety of formulations or topical nuances...
								\begin{compactitem}
									\item BDAG lists here an immense number of examples and ways of using this term.
								\end{compactitem}
							\item[1.a.$\gamma$] of an individual declaration or remark: \textit{assertion, declaration, speech}
							\item[1.a.$\delta$] the plural (\gr{οἱ}) \gr{λόγοι} is used, on the one hand, of words uttered on various occasions, of speeches or instruction given here and there by humans or transcendent beings...Also of words and expressions that form a unity, whether it be connected discourse...a conversation, or parts of one and the same teaching, or expositions on the same subject...
							\item[1.a.$\epsilon$] \textit{the subject} under discussion, \textit{matter, thing} generally
							\item[1.b] of literary or oratorical productions: of the separate books of a work...\textit{treatise}...
							\item[2] \textbf{\textit{computation, reckoning}}
							\item[2.a] a formal accounting, especially of one's actions, and frequently with figurative extension of commercial terminology \textit{account, accounts, reckoning}...Of a \textit{ledger heading}...\hyperref[ROMANS-9-6]{Romans 9:6} (the point is that God's `list' of Israelites is accurate; on \gr{ἐκπίπτω} in the sense `is not deficient'...\hyperref[ROMANS-9-9]{vs. 9} (the `count' is subsumed by metonymy in divine promise); \hyperref[GALATIANS-5-14]{Galatians 5:14} (all moral obligations come under one `entry': `you shall love your neighbor as yourself'...the contexts of these three passages suggest strong probability for commercial associations; for another view see \#1.a.$\beta$.
							\item[2.b] \textit{settlement} (\textit{of an account})
							\item]2.c] \textit{reflection, respect, regard}
							\item[2.d] reason for or cause of something, \textit{reason, ground, motive}
							\item[2.e] \gr{πρὸς ὸν ἡμῖν ὁ λόγος (ἐστίν)} \textit{with whom we have to do} (i.e. \textit{to reckon})
							\item[3] \textbf{the independent personified expression of God, \textit{the Logos}}. Our literature shows traces of a way of thinking that was widespread in contemporary syncretism, as well as in Jewish wisdom literature and Philo, the most prominent feature of which is the concept of the Logos, the independent, personified `Word' (of God)...\hyperref[JOHN-1-1]{John 1:1}, \hyperref[JOHN-1-14]{14}...It is the distinctive teaching of the Fourth Gospel that this divine `Word' took on human form in a historical person, that is, in Jesus...Equating a divinity with an abstraction that she personifies: Artem. 5, 18 \gr{φρόνησις εἶναι νομίζεται ἡ θεός} [Athena]). Cp. \hyperref[1JOHN-1-1]{1 John 1:1}; \hyperref[REVELATIONNT-19-13]{Revelation 19:13}. \gr{εἷς θεός ἐστιν, ὁ φανερώσας ἑαυτὸν διὰ Ἰ. Χριστοῦ τοῦ υἱοῦ αὐτοῦ, ὅς ἐστιν αὐτοῦ λόγος, ἀπὸ σιγῆς προελθών} \textit{there is one God, who has revealed himself through Jesus Christ his Son, who is his `Word' proceeding from silence} (i.e., without an oral pronouncement: in a transcendent manner) IMg 8:2 (s. \gr{σιγή}). The Lord as \gr{νόμος κ. λόγος} PtK 1. Cp. Dg 11:2, 3, 7, 8; 12:9.---HClavier, TManson memorial vol., '59, 81--93: the Alexandrian eternal \gr{λόγος} is also implied in \hyperref[HEBREWS-4-12]{Hebrews 4:12}; \hyperref[HEBREWS-13-7]{13:7}.
								% PtK = Petruskerygma (Preaching of Peter)— Petruskerygma (Preaching of Peter): EKlostermann, Apokrypha I (Kl. T. 3) 1933 pp. 13–16 (by p. and line)
								% Dg = Letter of Diognetus, attributed to an unidentifiable apologist; Diognetus. Bihlmeyer 141–49 (Bihlmeyer = KB., Die Apostolischen Väter: Neubearbeitung der Funkschen Ausgabe, Pt I, 2d ed. 1956)
						\end{compactitem}
					% \item[CGL , PDF ] \textbf{}
						% \begin{compactitem}
							% \item 
						% \end{compactitem}
					\item[LSJ , PDF ] \textbf{}
						\begin{compactitem}
							\item[I] \textit{computation, reckoning}
							\item[I.1.a] \textit{account} of money handled
							\item[I.1.b] \textit{public accounts}, i.e. branch of \textit{treasury}
							\item[I.2] generally, \textit{account, reckoning}
							\item[I.3] \textit{measure, tale}
							\item[I.4] \textit{esteem, consideration, value} put on a person or thing
							\item[II] \textit{relation, correspondence, proportion}
							\item[II.1] generally, \gr{ὑπερτερίης λ}. relation (of gold to lead)...
							\item[II.2] Mathematics, \textit{ratio, proportion}
							\item[II.3] Grammar, \textit{analogy, rule}
							\item[III] \textit{explanation}
							\item[III.1.a] \textit{plea, pretext, ground}
							\item[III.1.b] \textit{plea, case}
							\item[III.2.a] \textit{statement of a theory, argument}
							\item[III.2.b] \gr{ὁ περὶ θεῶν λ}., title of a discourse by Protagoras...
							\item[III.2.c] In logic, \textit{proposition}, whether as premise or conclusion
							\item[III.2.d] \textit{rule, principle, law}, as embodying the result of \gr{λογισμός}
						\end{compactitem}
					\item[TDNT] \href{https://drive.google.com/file/d/1pbXBjj7Em3K7JqcauiaXpz-RphcX6cdj/view?usp=sharing}{4:69--136}.
					\item[Dan Graham, ``Once More unto the Stream,'' 2013, 318--319] In the same way, it appears that for Heraclitus the whole cosmos is everlasting precisely because parts of it are constantly undergoing material changes (B30). Thus we seem to be justified in seeing the flux doctrine as having application to the cosmos as well as to rivers and people. In Heraclitus’ world, relatively stable structures supervene on constantly changing matter. B 12 is the river fragment, Heraclitus’ one and only statement of the flux. The sage of Ephesus believes in the flux of matter, but not a flux that erodes all landmarks and makes knowledge and communication impossible. Rather, his flux brings stability to rivers, life to animals, coherent experience to men, and everlasting existence to the cosmos itself. Or rather, there are regularities that govern flux, that order the exchanges of opposite characteristics. While the changes are continuous, the regularities balance them out and ensure a higher-level order in the world. Like Adam Smith’s Invisible Hand that maintains the balance of market forces in an economy, Heraclitus’ scheme maintains a balance of material forces. The mysterious entity he calls the Logos can be seen as just the expression of systematic order in nature. No less real than the constant exchange of hot and cold, wet and dry, earth, water, and fire is the constant replacement of what is lost by a new supply, drawn from what is lost of the contrary stuff. The world needs constant flux and interchange at the level of material, on which a hidden law operates to maintain balance. The cosmos is an ever-living fire precisely because it lives in its own transformation of opposites; it stays the same as it changes. There is one way that Heraclitus’ theory is profoundly different from that of all other early philosophers of nature. Whereas they attempt to explain change in terms of constant realities and their modifications, Heraclitus attempts to explain constancy in terms of ongoing change. Whereas in some metaphysical sense they are all philosophers of stasis, for whom (especially from Parmenides on) stability is fundamental and change problematic, Heraclitus is a philosopher of process, for whom change is fundamental and stability problematic. But if he is a process philosopher, he is not a nihilist or an epistemological or ontological anarchist. His world is orderly through and through, the forces balanced, the elements proportionate, the cosmos stable. Like a river.
				\end{compactdesc}
				
				% \begin{tabular}{ll}
					% Total occurrences in the NT & 
				% \end{tabular}
				
				\begin{compactdesc}
					\item[Some NT uses] \textbf{}
						\begin{compactdesc}
							\item[{\hyperref[JOHN-5-38]{John 5:38}}]
						\end{compactdesc}
				\end{compactdesc}
		
	% % -----
	% \subsubsection{Book of Mormon} % *BOM
	% % -----
		% \label{WORD-BOM}
	
		% \begin{tabular}{ll}
			% Total occurrences in the BOM & 
		% \end{tabular}
		
		% \begin{compactdesc}
			% \item[]
		% \end{compactdesc}
		
	% % -----
	% \subsubsection{Doctrine and Covenants} % *DC
	% % -----
		% \label{WORD-DC}
	
		% \begin{tabular}{ll}
			% Total occurrences in the DC & 
		% \end{tabular}
		
		% \begin{compactdesc}
			% \item[]
		% \end{compactdesc}
		
	% % -----
	% \subsubsection{Pearl of Great Price} % *PGP
	% % -----
		% \label{WORD-PGP}
	
		% \begin{tabular}{ll}
			% Total occurrences in the PGP & 
		% \end{tabular}
		
		% \begin{compactdesc}
			% \item[]
		% \end{compactdesc}
		
% ----------------------
\subsection{Key Phrases} % *KEYPHRASES *PHRASES
% ----------------------
	\label{WORD-PHRASES}

	\begin{compactdesc}
	
		\item[word(s) of eternal life] \textbf{3} | \hyperref[JOHN-6-68]{John 6:68}; \hyperref[DC84-43]{DC 84:43}; \hyperref[MOSES-6-59]{Moses 6:59} % "word* of* et* li*"
	
		\item[word(s) of life] \textbf{3} | \hyperref[PHILIPPIANS-2-16]{Philippians 2:16}; \hyperref[1JOHN-1-1]{1 John 1:1}; \hyperref[DC84-85]{DC 84:85} % "word* of* li*"
	
		\item[words of my mouth] \textbf{10} | \hyperref[DEUTERONOMY-32-1]{Deuteronomy 32:1}; \hyperref[PSALMS-19-14]{Psalms 19:14}; \hyperref[PSALMS-54-2]{54:2}; \hyperref[PSALMS-78-1]{78:1}; Proverbs 4:5; 5:7; 7:24; 8:8; Hosea 6:5; \hyperref[ALMA-7-1]{Alma 7:1}
			\begin{compactdesc}
				% \item[Bible anchor verses] \phantom{.}
					% \begin{compactdesc}
						% \item[XXX] 
					% \end{compactdesc}
				\item[Commentary] See discussion in \hyperref[ENDOWMENT-words-mouth]{this note} in the endowment ceremony.
			\end{compactdesc}
			
		\item[word of my power/power of my word] \textbf{17} | \hyperref[LUKE-4-32]{Luke 4:32}, \hyperref[HEBREWS-1-3]{Hebrews 1:3}, \hyperref[HEBREWS-4-12]{Hebrews 4:12}, \hyperref[1NEPHI-17-46]{1 Nephi 17:46}, \hyperref[2NEPHI-1-26]{2 Nephi 1:26}, \hyperref[JACOB-4-9]{Jacob 4:9}, \hyperref[ALMA-5-5]{Alma 5:5}, \hyperref[ALMA-26-13]{Alma 26:13}, \hyperref[4NEPHI-1-30]{4 Nephi 1:30}, \hyperref[MORMON-8-24]{Mormon 8:24} (2), \hyperref[MORMON-9-17]{Mormon 9:17} (2), \hyperref[DC11-21]{DC 11:21}, \hyperref[MOSES-1-32]{Moses 1:32}, \hyperref[MOSES-1-35]{Moses 1:35}, \hyperref[MOSES-2-5]{Moses 2:5}
			% \begin{compactdesc}
				% \item[Bible anchor verses] \phantom{.}
					% \begin{compactdesc}
						% \item[XXX] 
					% \end{compactdesc}
				% \item[Commentary] 
			% \end{compactdesc}
			
	\end{compactdesc}
	
% % ----------------------
% \subsection{Bible Dictionaries} % *DICTIONARIES
% % ----------------------
	% \label{WORD-BD}

% % ----------------------
% \subsection{Restoration Usage} % *RESTORATION
% % ----------------------
	% \label{WORD-RESTORATION}

	% % -----
	% \subsubsection{Joseph Smith} % *JS
	% % -----
		% \label{WORD-JS}
	
		% \begin{compactdesc}
			% \item[{\hyperref[KEY]{Date/Source}}]
		% \end{compactdesc}
		
	% % -----
	% \subsubsection{Temple Ordinances}
	% % -----
		% \label{WORD-TEMPLE}
	
		% \begin{compactdesc}
			% \item[{\hyperref[KEY]{Date/Source}}]
		% \end{compactdesc}
	
	% % -----
	% \subsubsection{Brigham Young} % *BY
	% % -----
		% \label{WORD-BY}
	
		% \begin{compactdesc}
			% \item[{\hyperref[KEY]{Date/Source}}]
		% \end{compactdesc}
	
% % ----------------------
% \subsection{Commentary and Studies} % *COMMENTARY *STUDIES
% % ----------------------
	% \label{WORD-COMMENTARY}

	% % -----
	% \subsubsection{Key Points}
	% % -----
		% \label{WORD-KEYS}
	
		% \begin{compactitem}
			% \item
		% \end{compactitem}
		
	% % -----
	% \subsubsection{Sample Study}
	% % -----
		% \label{WORD-study}